% =====================================================================
%  chunks/concl.tex —— 结语 Concluding Remarks（原书 p.273，PDF 280）
% =====================================================================
% [原文 p.273]
\chapter*{结语}
\addcontentsline{toc}{chapter}{结语}
\markboth{结语}{结语}

在这篇关于拓扑学的观念与方法的综述中，远非所有课题都已涉及。关于 $3$ 维拓扑学，以及由 Thurston 等人发展的 Kleinian 群理论，这里谈得很少；关于 Donaldson、Freedman 等人带来的 $4$ 维拓扑学的进展，也谈得很少。我们同样没有综述叶状结构的定性理论——尤其是 Reeb、Haefliger 与 Novikov 关于解析叶状结构以及具紧叶的叶状结构的定理。此外还有：建立开流形上叶状结构及其他几何结构存在性的技术（Gromov）；奇数维球面上余维一叶状结构的存在性（Lawson、Tamura 等人）；任意维紧流形的相应问题（Thurston）；叶状结构的示性类理论（Bott、Godbillon、Vey、Bernshtein 等人）；叶状结构的分类理论（Haefliger 等人）。最后是流形上向量场 Lie 代数的上同调理论（I.\,M. Gel'fand、Fuks）、复超曲面的奇点理论（Milnor、Brieskorn）与实代数曲线及曲面（Gudkov、Arnol'd、Rohlin、Kharlamov、Kirilov、Viro）。

我们也从综述中略去了拓扑学的如下领域：流形嵌入理论与高维纽结理论（由若干人创立，始于 Whitney，包括 Haefliger、Stallings 与 Levine），以及「映射与函数的典型奇点」理论（Whitney、Pontryagin、Thom、Boardman、Mather、Arnol'd）。

本综述的作者\emph{有意地}略去了对一般范畴性的、抽象性的拓扑结果的讨论，尽管就拓扑概念的内在逻辑而言，这些结果无论多么有用、甚至多么必要。

综述中还缺了关于一般空间同调理论的概述，特别是 $\R^n$ 的子集的同调理论；这一理论在 20 世纪 20 年代由若干作者发展起来，其中包括 P.\,S. Alexandrov、Čech、Pontryagin、Kolmogorov、Steenrod、Chogoshvili、Sitnikov、Milnor 等人。

还缺了关于具有各种微分几何性质的流形的拓扑性质的讨论（整体几何）。

最后，拓扑学近几十年来对实际物理问题的实质性应用——这些应用已改变了现代数学物理的整套工具——未能加以讨论。我们希望这一缺憾将在本丛书的其他各卷中得到弥补。

% [Tutor Note: 上文的 Bernshtein（原书如此拼写）即 I. N. Bernshtein，今多写作 Bernstein；Kirilov 原书如此，今多写作 Kirillov。两人名的拼写均照原书，不作径改。]
% [Tutor Note: 这份「结语」其实是一张缺项清单，把 1986 年（英译本 1996 年）时作者认为本卷未及展开的方向逐一点名。读它可以当索引：凡这里列出的方向，正文中至多一笔带过，不要指望在本书里找到系统论述。这也解释了附录为何专谈 3 维流形与纽结——正是作者认为正文最欠缺的部分之一。]

% CHUNK concl DONE
